\documentclass[12pt]{article}

\usepackage[margin=1in]{geometry}
\usepackage{setspace}
\usepackage{amsmath, amssymb, amsthm}
\usepackage{graphicx}
\usepackage[superscript,biblabel]{cite}
\usepackage{float}
\usepackage[hidelinks]{hyperref}
\usepackage{caption}
\usepackage{parskip}

\singlespacing
\setlength{\parindent}{1pt}

\title{APMA 0365 Final Project}
\author{Finn Quested, Matthew Williamson, Matthew Purcell, Gabin Rwigema, Ines Saltiel}
\date{May 2026}

\begin{document}
\begin{titlepage}
    \centering
    \vspace*{\fill}
    {\Huge \textbf{Extending the Ross Malaria Model to Reaction-Diffusion PDEs: Spatial Dynamics, Diffusion Thresholds, and Mosquito Control} \\[0.5cm]
    {\Large Finn Quested, Matthew Williamson, Matthew Purcell, Gabin Rwigema, Ines Saltiel} \\[0.4cm]
    May 6, 2026
    \vfill
\end{titlepage}
\newpage
\tableofcontents
\newpage
\section{Introduction and Context}

Malaria is one of the most significant infectious diseases in the world. The World Health Organization estimated 249 million cases and over 600,000 deaths in 2022 alone, mainly affecting children under five in sub-Saharan Africa \cite{Qu2023}. 

A key challenge in controlling malaria is that transmission is highly spatially heterogeneous. \textit{Anopheles} mosquitoes, the primary carrier of the parasite, require water to breed and thrive in humid conditions, so transmission is dependent on local ecology. People traveling between regions can also introduce the disease into new areas or encounter it in new places. Thus, malaria risk is unevenly distributed across space and depends on interactions between human and mosquito populations.

The transmission of malaria is bidirectional. An infectious mosquito injects parasites into a human host, where they mature in the liver before eventually producing a transmissible form. Then, when another mosquito bites this person, it ingests these parasites, before being passed to the next human host. This is what motivates the classical two-population structure of the Ross model, which will be discussed in greater depth later.

\section{Introducing the ODE model}

\subsection{The Classical Ross Model}

We first consider the classical ODE Ross-type model for malaria transmission between humans and mosquitoes. Ronald Ross presented a system of differential equations to describe the transmission of malaria in a population, focusing on the role of mosquitoes \cite{Ross1911}. Based on the Susceptible, Infected, Recovered (SIR) framework, humans and the mosquitoes are divided into “susceptible” and “infected” compartments. The infected human population becomes susceptible again after the incubation time elapses, whereas the mosquito population stays infected. Thus, this is an “SIS” model and “SI” model for humans and mosquitoes, respectively. 

The spread is modeled by the compartment model:
\begin{figure}[H]
    \centering
    \includegraphics[width=0.8\linewidth]{Ross Model Representation.png}
    \caption{Compartment model for the Ross-type malaria system, with total human and mosquito populations held constant ($S_h + I_h = 1$, $S_m + I_m = 1$). Arrow labels are per-capita transition rates; the total flow along an arrow equals the rate times the size of the source compartment (e.g., the $S_h \to I_h$ flow is $abm \cdot I_m \cdot (1 - I_h)$). The $I_m \to S_m$ arrow represents death of infected mosquitoes and their replacement by new susceptibles under the constant-population assumption, so the mosquito compartment is biologically SI rather than SIS.}
    \label{fig:compartment_ode}
\end{figure}

The system of differential equations presented by the Ross model is described by: 

\[
\begin{cases}
{{dS_h}\over{dt}} = rI_h(t) - abmI_m(t)(1-I_h(t)),\\
{{dI_h}\over{dt}} = abmI_m(t)(1-I_h(t))-rI_h(t),\\
{{dS_m}\over{dt}} = \mu I_m(t) - aqI_h(1-I_m(t)),\\
{{dI_m}\over{dt}} = aqI_h(t)(1-I_m(t))-\mu I_m(t).
\end{cases}
\]
Where 
\begin{itemize}
    \item $I_h(t)$: fraction of infected humans
    \item $I_m(t)$: fraction of infected mosquitoes
    \item $S_h(t)$: fraction of susceptible humans at time t
    \item $S_m(t)$: fraction of susceptible mosquito at time t
    \item $a$: bite rate per mosquito
    \item $b$: probability an infectious mosquito infects a human
    \item $q$: probability a mosquito becomes infectious after biting an infected human
    \item $m$: mosquito-to-human ratio
    \item $r$: human recovery rate
    \item $\mu$: mosquito death rate,
\end{itemize}

We assume the mosquito population is at constant size, with the death rate $\mu$ matched by an equal birth rate of susceptible mosquitoes. The term $-\mu I_m$ in the $I_m$ equation reflects loss of infected mosquitoes through this constant-population mortality, with replacements entering the susceptible pool. We also assume humans and mosquitoes become infectious instantly upon transmission.

So,
\[
S_h(t)+I_h(t) = 1, \quad S_m(t) + I_m(t) = 1.
\]
Giving initial conditions:
\[
\begin{cases}
S_h(0) + I_h(0) = 1\\
S_m(0) + I_m(0) = 1.
\end{cases}
\]

 The strength of this cycle determines whether malaria dies out (infection levels approach 0), or persists at an endemic equilibrium.

Since:
\[
S_h(t) = 1 - I_h(t), \quad S_m(t) = 1 - I_m(t),
\]
The full system can be reduced to
\[
\begin{cases}
\frac{dI_h}{dt} = abm I_m(t)(1 - I_h(t)) - r I_h(t), \\
\frac{dI_m}{dt} = aq I_h(t)(1 - I_m(t)) - \mu I_m(t).
\end{cases}
\]

We nondimensionalize time as \(\tau = rt\) and introduce the dimensionless parameters
\[
\alpha = \frac{abm}{r}, \qquad \delta = \frac{aq}{\mu}, \qquad \varepsilon = \frac{r}{\mu}.
\]
Here \(\alpha\) measures the rate at which infectious mosquitoes produce new human infections relative to human recovery, \(\delta\) measures the rate at which infected humans produce new mosquito infections relative to mosquito mortality, and \(\varepsilon\) is the ratio of human to mosquito time scales. Since \(\varepsilon=r/\mu\), small \(\varepsilon\) means mosquitoes evolve on a faster time scale than humans.

This yields the nondimensionalized system:
\[
\begin{cases}
\frac{dI_h}{d\tau} = \alpha I_m(1 - I_h) - I_h, \\
\frac{dI_m}{d\tau} = \frac{\delta}{\varepsilon} I_h(1 - I_m) - \frac{1}{\varepsilon} I_m.
\end{cases}
\]

Both \(I_h(\tau)\) and \(I_m(\tau)\) represent fractions of infected populations and must lie in \([0,1]\). For the uncontrolled nondimensional Ross model, the disease-free equilibrium is
\[
(I_h,I_m)=(0,0),
\]
and the basic reproduction number is
\[
\mathcal{R}_0=\sqrt{\alpha\delta}.
\]
When \(\mathcal{R}_0<1\), small infections decay near the disease-free equilibrium; when \(\mathcal{R}_0>1\), small infections grow.

\subsection{Limitation of the ODE Model}

Alemu et al. show that malaria attacks are not evenly distributed in space and time. In highland areas with low endemicity, malaria transmission is extremely variable, and infection risk for individuals is unevenly distributed \cite{Alemu2014}. Therefore, while the Ross model captures the temporal dynamics of malaria transmission and predicts the conditions for disease persistence ($\mathcal{R}_0 > 1$), it assumes a well-mixed population.

This assumption leads to several critical limitations:

\begin{itemize}
    \item \textbf{No spatial heterogeneity}: The model cannot distinguish between a localized outbreak and a region-wide epidemic. An outbreak affecting 10\% of the population looks the same, concentrated in a single neighborhood or spread evenly across an entire region.
    
    \item \textbf{No spatial hotspots}: Areas with high mosquito density near water, or regions with poor access to healthcare, cannot be identified or modeled separately.
    
    \item \textbf{No spreading dynamics}: The model cannot describe how malaria spreads from an endemic region into surrounding areas, or the speed in which the disease is invading a region.  
    
    \item \textbf{No localized interventions}: Inability to evaluate control strategies that target specific outbreak areas, including indoor residual spraying in high-transmission zones.
\end{itemize}

A surveillance study in northwest Ethiopia identified 56 statistically significant spatial clusters of malaria cases in high-altitude villages
(Figure~\ref{fig:ethiopia_spatial}). The spatial scan statistics revealed
that 75.7\% of cases occurred in villages above 2000 meters in elevation,
concentrated in specific geographic clusters rather than distributed evenly
across the region~\cite{Alemu2014}.

This spatial heterogeneity fundamentally violates the well-mixed assumption
of the ODE model. Consider two scenarios that produce identical predictions
in the ODE framework:
\begin{itemize}
    \item \textbf{Scenario A:} 645 malaria cases distributed uniformly
    across the 159{,}091 people in the district, giving the same low
    infection density everywhere.

    \item \textbf{Scenario B:} The same 645 cases, but with 488 of them
    (75.7\%) concentrated in a handful of high-altitude villages and the
    remaining 157 scattered across the rest of the district.
\end{itemize}

Both scenarios produce the same total infected fraction, so the ODE model cannot tell them apart. The Ethiopian data clearly correspond to Scenario B, where infection is spatially clustered rather than uniformly mixed. This motivates the development of spatially explicit PDE models that distinguish between these cases.

\begin{figure}[H]
    \centering
    \includegraphics[width=0.85\linewidth]{ehtiopiafig2.png}
    \caption{Spatial clustering of malaria cases at the administrative unit level in northwest Ethiopia. Cases are highly concentrated in specific regions  (shown in red/orange) rather than uniformly distributed. The color scale represents log-likelihood ratios (LLR) indicating the strength of clustering.  Out of 645 malaria cases, 75.7\% occurred in villages above 2,000m 
    elevation. The ODE model cannot capture this spatial heterogeneity. 
    Taken from Alemu et al. (2014).}
    \label{fig:ethiopia_spatial}
\end{figure}
\newpage

\section{From ODE to PDE}

\subsection{Adding Space}

To address these limitations and capture the spatial nature of malaria transmission, we extend the ODE model to include spatial variation. Instead of tracking only the total infected fractions \(I_h(\tau)\) and \(I_m(\tau)\), we track local infected fractions \(I_h(x,\tau)\) and \(I_m(x,\tau)\), where \(x\) represents position and \(\tau\) is nondimensional time. This leads to a system of partial differential equations.

We introduce a one-dimensional spatial domain \([0,L]\), representing a geographic transect or region of interest. In the context of the Ethiopian study, this spatial domain can be interpreted as a transect across villages with varying elevation and transmission intensity. The ODE variables are replaced as follows:
\begin{align}
I_h(\tau) &\longrightarrow I_h(x,\tau)
\quad \text{(local infected human fraction at location \(x\) and time \(\tau\))}, \\
I_m(\tau) &\longrightarrow I_m(x,\tau)
\quad \text{(local infected mosquito fraction at location \(x\) and time \(\tau\))}.
\end{align}

Similarly, the susceptible fractions become spatially dependent:
\begin{align}
S_h(\tau) &\longrightarrow S_h(x,\tau), \\
S_m(\tau) &\longrightarrow S_m(x,\tau).
\end{align}

\subsection{Interpretation of Spatial Fractions}

The function \(I_h(x,\tau)\) represents the local fraction of infected humans at position \(x\) and time \(\tau\). For example, if \(I_h(5,\tau)=0.20\), then approximately 20\% of the local human population near position \(x=5\) is infected at that time. This type of spatial variation reflects patterns observed in the Ethiopian data, where infection levels varied significantly between clustered high-risk villages and surrounding areas \cite{Alemu2014}.

To summarize infection over the whole domain, we can define the spatial average
\[
\overline{I}_h(\tau)=\frac{1}{L}\int_0^L I_h(x,\tau)\,dx.
\]
This quantity is the PDE analogue of the infected fraction tracked by the ODE model. The shift from \(I_h(\tau)\) to \(I_h(x,\tau)\) changes the modeling perspective: the PDE model predicts where infections are concentrated and how spatial patterns evolve over time.

\subsection{The Reaction-Diffusion PDE System}

To derive the PDE system, we account for two processes:

\begin{enumerate}
    \item Reaction terms: the same infection and recovery dynamics from the ODE model, now acting locally at each point $x$.
    \item Diffusion terms: the spatial movement of infected humans and mosquitoes.
\end{enumerate}

Combining the local Ross reaction terms with the spatial  diffusion creates a reaction-diffusion system. The diffusion mechanism is derived in Section 4, and the complete PDE system is assembled in Section 5. 

\subsection{Comparing the ODE and PDE analytically}

The fundamental difference between the ODE and PDE models lies in what each can describe:

\textbf{ODE Model Capabilities:}
\begin{itemize}
    \item Predicts the total number of infected humans $I_h(t)$ and infected mosquitoes $I_m(t)$ over time
    \item Determines whether malaria persists ($\mathcal{R}_0 > 1$) or dies out ($\mathcal{R}_0 \leq 1$)  
    \item Provides global outbreak dynamics without spatial information
\end{itemize}

\textbf{PDE Model Capabilities:}
\begin{itemize}
    \item Predicts \emph{where} infection is concentrated: $I_h(x,t)$ shows which locations have high infection density
    \item Describes \emph{how} infection spreads spatially over time
    \item Can model localized outbreaks that remain confined to certain regions
    \item Enables analysis of spatially targeted interventions
\end{itemize}

\subsection{Low versus High Mosquito Diffusion}

The mosquito diffusion coefficient $D_m$ controls how quickly infected mosquitoes spread through space:

\begin{itemize}
    \item \textbf{Low $D_m$} (mosquitoes stay near breeding sites): Infection remains localized. Even if $\mathcal{R}_0 > 1$ locally, malaria may not spread to distant regions. 
    
    \item \textbf{High $D_m$} (mosquitoes disperse widely): Infected mosquitoes carry malaria far from the initial outbreak location. This leads to wider spatial spread of human infections. 
\end{itemize}

The human diffusion coefficient $D_h$ plays a similar role, though human movement is typically more directed (travel between villages) than mosquito movement.

\subsection{No Control versus Mosquito Control}

Section 8 introduces mosquito control as an additional removal term in the infected mosquito equation. The PDE framework allows us to predict how control effectiveness depends on spatial factors:

\begin{itemize}
    \item \textbf{Without control \((c=0)\):} Infected mosquitoes are removed at the natural nondimensional rate \(1/\varepsilon\). If \(\mathcal{R}_0>1\), malaria can persist and may spread spatially depending on \(D_m\).

    \item \textbf{With control \((c>0)\):} The effective infected-mosquito removal rate becomes
    \[
    \frac{1}{\varepsilon}+c.
    \]
    This lowers the controlled reproduction number and can shift the system from outbreak growth toward decay near the disease-free equilibrium. Spatially, control can reduce transmission in the domain while diffusion still redistributes infection across space.
\end{itemize}

\section{Derivation of the Diffusion Term}

The previous section introduced the idea that infected humans and infected mosquitoes should be modeled as spatial densities rather than only as time-dependent populations. We now derive the diffusion term that accounts for movement through space. The goal is to explain why the PDE model should contain second spatial derivatives, such as
\[
D_h \frac{\partial^2 I_h}{\partial x^2}
\qquad \text{and} \qquad
D_m \frac{\partial^2 I_m}{\partial x^2}.
\]
This is the same mathematical structure as the heat equation studied in class, except that here the quantity spreading through space is infection density rather than heat.

Let \(I(x,t)\) represent a general infection fraction at position \(x\) and time \(t\). This could represent either \(I_h(x,t)\), the infected human fraction, or \(I_m(x,t)\), the infected mosquito fraction. To derive the spatial part of the PDE, we begin with a conservation law on a small interval \([a,b]\subset [0,L]\). The total amount of infection in this interval is
\[
\int_a^b I(x,t)\,dx.
\]
This quantity can change for two reasons. First, infected individuals may move across the endpoints of the interval. Second, infection may be created or removed locally through transmission, recovery, or mosquito death. Therefore, we write
\[
\frac{d}{dt}\int_a^b I(x,t)\,dx
=
J(a,t)-J(b,t)+\int_a^b R(I)(x,t)\,dx,
\]
where \(J(x,t)\) is the flux of infected fraction and \(R(I)\) represents the local reaction terms. The sign convention is that \(J(a,t)\) is flux entering through the left endpoint, while \(J(b,t)\) is flux leaving through the right endpoint.

Using the Fundamental Theorem of Calculus, the flux contribution can be rewritten as
\[
J(a,t)-J(b,t)
=
-\int_a^b \frac{\partial J}{\partial x}(x,t)\,dx.
\]
Substituting this into the balance law gives
\[
\int_a^b \frac{\partial I}{\partial t}(x,t)\,dx
=
\int_a^b \left(-\frac{\partial J}{\partial x}(x,t)+R(I)(x,t)\right)\,dx.
\]
Since this identity holds for any interval \([a,b]\), the integrands must agree, so
\[
\frac{\partial I}{\partial t}
=
-\frac{\partial J}{\partial x}+R(I).
\]
This is the local conservation form of the model.

To close the equation, we need a rule for the flux \(J\). We assume Fickian diffusion, meaning that infected density tends to move from regions where it is high toward neighboring regions where it is lower. Mathematically, this is written as
\[
J=-D\frac{\partial I}{\partial x},
\]
where \(D>0\) is the diffusion coefficient. If \(I_x>0\), then the density is increasing as \(x\) increases, so the flux points to the left. If \(I_x<0\), then the density is decreasing as \(x\) increases, so the flux points to the right. This matches the idea that diffusion smooths out sharp spatial differences.

Substituting Fick's law into the conservation equation gives
\[
\frac{\partial I}{\partial t}
=
-\frac{\partial}{\partial x}\left(-D\frac{\partial I}{\partial x}\right)+R(I).
\]
Assuming \(D\) is constant in space, this simplifies to
\[
\frac{\partial I}{\partial t}
=
D\frac{\partial^2 I}{\partial x^2}+R(I).
\]
Thus, the diffusion term in one spatial dimension is \(D I_{xx}\). In a higher-dimensional spatial region, the same idea would give
\[
\frac{\partial I}{\partial t}
=
D\Delta I+R(I),
\]
where \(\Delta\) is the Laplacian.

This derivation connects directly to the heat equation from class:
\[
u_t=Du_{xx}.
\]
For the heat equation, \(u(x,t)\) represents temperature or particle concentration. In this malaria model, \(I_h(x,t)\) and \(I_m(x,t)\) represent infected population fractions at each location. The interpretation is different, but the diffusion mechanism is the same: high-concentration regions spread into nearby lower-concentration regions over time.

\section{Reaction-Diffusion Malaria Model}

We now combine the local malaria transmission dynamics from the Ross model with the diffusion mechanism derived above. The model tracks two spatially dependent quantities:
\[
I_h(x,t) = \text{fraction of infected humans at location } x \text{ and time } t,
\]
\[
I_m(x,t) = \text{fraction of infected mosquitoes at location } x \text{ and time } t.
\]
The spatial domain is the one-dimensional interval \(x\in[0,L]\), which represents a simplified geographic region or transect. The model assumes that susceptible fractions can be written as
\[
S_h(x,t)=1-I_h(x,t),
\qquad
S_m(x,t)=1-I_m(x,t),
\]
so the state variables remain the infected fractions.

The reaction-diffusion system is
\begin{equation}
\frac{\partial I_h}{\partial t}
=
D_h\frac{\partial^2 I_h}{\partial x^2}
+
\alpha I_m(1-I_h)
-
I_h,
\label{eq:human_rd}
\end{equation}
\begin{equation}
\frac{\partial I_m}{\partial t}
=
D_m\frac{\partial^2 I_m}{\partial x^2}
+
\frac{\delta}{\varepsilon} I_h(1-I_m)
-
\frac{1}{\varepsilon} I_m.
\label{eq:mosquito_rd}
\end{equation}
Here \(D_h\) and \(D_m\) are diffusion coefficients for infected humans and infected mosquitoes. 

Equation \eqref{eq:human_rd} describes the infected human fraction. The term
\[
D_h\frac{\partial^2 I_h}{\partial x^2}
\]
models spatial movement of infected humans. If infected humans are concentrated in one part of the domain, this term causes that concentration to spread into neighboring regions. The term
\[
\alpha I_m(1-I_h)
\]
models new human infections caused by infected mosquitoes, where \(\alpha\) represents the human transmission rate. It depends on the local infected mosquito fraction \(I_m(x,t)\) and the local susceptible human fraction \(1-I_h(x,t)\). Finally,
\[
-I_h
\]
models human recovery. Because time has been nondimensionalized by the human recovery rate, the recovery coefficient is \(1\).

Equation \eqref{eq:mosquito_rd} describes the infected mosquito fraction. The term
\[
D_m\frac{\partial^2 I_m}{\partial x^2}
\]
models spatial movement of infected mosquitoes. This coefficient is especially important biologically because mosquitoes may remain near breeding sites or disperse across a larger area depending on environmental conditions. The term
\[
\frac{\delta}{\varepsilon} I_h(1-I_m)
\]
models new mosquito infections caused by mosquitoes biting infected humans, scaled by the time-scale parameter \(\varepsilon\). It depends on the local infected human fraction \(I_h(x,t)\) and the local susceptible mosquito fraction \(1-I_m(x,t)\). The final term
\[
-\frac{1}{\varepsilon} I_m
\]
models infected mosquito death, where \(1/\varepsilon\) is the mosquito mortality rate relative to the human recovery rate.

The main modeling idea is that the reaction terms act locally, while the diffusion terms connect nearby locations. If there were no diffusion, then each point \(x\) would behave like its own local ODE model. If there were no reaction terms, then the system would behave like two separate heat equations, simply smoothing the initial infected human and mosquito distributions. The reaction-diffusion system combines these two effects: malaria can grow or decay locally through transmission and removal, while also spreading spatially through movement.

\begin{table}[h]
\centering
\begin{tabular}{c|p{0.72\textwidth}}
\hline
\textbf{Symbol} & \textbf{Meaning} \\
\hline
\(x\) & Spatial position in the one-dimensional domain \([0,L]\). \\
\(t\) & Nondimensional time variable. \\
\(L\) & Length of the spatial domain. \\
\(I_h(x,t)\) & Fraction of infected humans at location \(x\) and time \(t\). \\
\(I_m(x,t)\) & Fraction of infected mosquitoes at location \(x\) and time \(t\). \\
\(1-I_h(x,t)\) & Fraction of susceptible humans at location \(x\) and time \(t\). \\
\(1-I_m(x,t)\) & Fraction of susceptible mosquitoes at location \(x\) and time \(t\). \\
\(D_h\) & Human diffusion coefficient, representing the strength of spatial movement of infected humans. \\
\(D_m\) & Mosquito diffusion coefficient, representing the strength of spatial movement of infected mosquitoes. \\
\(\alpha\) & Human transmission parameter. Controls the rate at which infected mosquitoes produce new human infections. \\
\(\delta\) & Mosquito transmission parameter (scaled by \(\varepsilon\)). Controls the rate at which infected humans produce new mosquito infections. \\
\(\varepsilon\) & Time-scale ratio between mosquito and human dynamics. Mosquito death rate is \(1/\varepsilon\) relative to human recovery. \\
\hline
\end{tabular}
\caption{Variables and parameters in the reaction-diffusion malaria model.}
\label{tab:rd_parameters}
\end{table}

It is also useful to write the same system in a compact form:
\begin{align}
\frac{\partial I_h}{\partial t} &=
D_h I_{h,xx}+\alpha I_m(1-I_h)-I_h\\
\frac{\partial I_m}{\partial t}
&=
D_m I_{m,xx}+\frac{\delta}{\varepsilon} I_h(1-I_m)-\frac{1}{\varepsilon} I_m.
\end{align}
This notation highlights the structure of the model. The \(I_{h,xx}\) and \(I_{m,xx}\) terms are the diffusion terms, while the remaining terms are the local reaction terms. Therefore, the model is a reaction-diffusion system: reaction changes infection levels at each location, and diffusion redistributes infection across space.


\section{Boundary and Initial Conditions}

To complete the PDE system, we must impose boundary and initial conditions.

For malaria, we consider a one-dimensional spatial domain $x \in [0, L]$, representing a geographic region. To model a closed system, we impose no-flux (Neumann) boundary conditions consistent with those introduced in class for the heat equation\cite{Sandstede2026}:
\begin{align}
\frac{\partial I_h}{\partial x}(0,t) &= \frac{\partial I_h}{\partial x}(L,t) = 0, \\
\frac{\partial I_m}{\partial x}(0,t) &= \frac{\partial I_m}{\partial x}(L,t) = 0.
\end{align}

These boundary conditions specify that neither infected individuals nor mosquitoes move in or out of the region. This lets the model preserve spatial variation without introducing external infection through the boundary. This captures differences between high- and low-density infection regions, which is a key limitation of the ODE model discussed earlier. 

We also state initial conditions:
\[
I_h(x,0)=I_{h,0}(x), \qquad I_m(x,0)=I_{m,0}(x),
\]
which describe how infection is distributed at the start of the PDE model. In the numerical simulations below, we choose these initial profiles to be localized Gaussian bumps centered at \(x_0\). This is consistent with the clustering behavior observed by Alemu et al., where malaria cases were concentrated in specific villages rather than evenly distributed across the district. The Gaussian profile gives a mathematical representation of an outbreak that begins near one location before spreading over time and space.

\section{Qualitative Analysis of the PDE System}

\subsection{Positivity}

In the model, \(I_h(x,t)\) and \(I_m(x,t)\) represent fractions of infected populations. As a result, biologically meaningful solutions should satisfy
\[
0 \leq I_h(x,t),\, I_m(x,t) \leq 1.
\]
The continuous reaction terms are consistent with this  region. At \(I_h=0\) and \(I_m=0\), the reaction terms point inward, while at \(I_h=1\) and \(I_m=1\), they do not push the solution above \(1\). Diffusion also preserves this through spatial smoothing. 

\subsection{Diffusion and the Application of the Maximum Principle}

The diffusion terms
\[
D_h \frac{\partial^2 I_h}{\partial x^2}, \quad D_m \frac{\partial^2 I_m}{\partial x^2}
\]
derived earlier describes how infection spreads across space. Regions with high infection density tend to disperse into neighboring regions with lower density. Over time, this produces a smoothing effect: sharp peaks flatten, and the distribution spreads across space. 

The maximum principle supports this smoothing behavior. Although the PDE system is nonlinear, it retains qualitative features of the heat equation discussed in class. The strong maximum principle suggests that diffusion prevents the creation of new interior maxima or minima, so extreme values of the solution are largely determined by the initial and boundary conditions.

In the context of the Ethiopian village-level data from Alemu et al., this diffusion term represents how infection initially concentrated in a cluster can begin to affect nearby locations through human and mosquito movement. 

This intuition will be made more precise in Section 9.2, where the solution is decomposed into spatial modes and higher-frequency components are shown to decay more rapidly.

\subsection{Total Infection and the Role of Diffusion}

To examine the effect of diffusion on total infection, we define
\[
H(t) = \int_0^L I_h(x,t)\,dx.
\]

Differentiating with respect to time gives
\[
\frac{dH}{dt} = \int_0^L \frac{\partial I_h}{\partial t}\,dx.
\]

Substituting the PDE into this expression, the diffusion term simplifies:
\[
\int_0^L \frac{\partial^2 I_h}{\partial x^2} \, dx
= \frac{\partial I_h}{\partial x}(L,t) - \frac{\partial I_h}{\partial x}(0,t) = 0,
\]
by the no-flux Neumann boundary conditions.

This shows that diffusion does not change total infection, but instead redistributes it across space. Any increase or decrease in total human infection is driven entirely by the reaction terms (transmission and recovery). An analogous argument applies to the mosquito population.

This behavior is illustrated in Figure~\ref{fig:noflux}, where diffusion redistributes infection within the domain while the no-flux boundary condition prevents any loss at the boundaries.

\begin{figure}[htbp]
\centering
\includegraphics[width=0.65\textwidth]{noflux_sketch.png}
\caption{Illustration of the boundary conditions on the domain $[0,L]$. The infection concentration shifts under diffusion, while the zero-gradient condition at 0 and L prevents outward flux, preserving total infection.}
\label{fig:noflux}
\end{figure}

\section{Representation of Mosquito Control Effect}

Because our PDE models includes space, it can also represent how interventions strategies affect transmission across a region. Mosquito control measures like insecticide spraying and bed nets reduce the number of infectious mosquitoes and weaken transmission. Alemu et al. shows that malaria risk can cluster geographically. We use this extension to answer how reducing mosquito transmission changes the spread of infection in the model.

We extend the PDE model to include mosquito control. This is done by modifying the mosquito equation to include an additional removal term $c$

\begin{cases}
\frac{\partial I_h}{\partial t} &=
D_h I_{h,xx}+\alpha I_m(1-I_h)-I_h \\
\frac{\partial I_m}{\partial t}
&= D_m \frac{\partial^2 I_m}{\partial x^2}
+ \frac{\delta}{\varepsilon} I_h(1 - I_m)
- \left(\frac{1}{\varepsilon} + c\right) I_m
\end{cases}

where $c > 0$ represents the strength of control measures such as insecticide spraying or bed nets. Mathematically, increasing $c$ lowers the number of infectious mosquitoes.

This extension affects only the reaction term and leaves the diffusion term unchanged. The overall level of infection is reduced, while the spatial spreading mechanism remains unchanged.

We model mosquito control as an added removal rate $c$ for infected mosquitoes. However, in the real world, control is not uniform and can vary across regions and over time. Our model captures the main effect of increased mosquito death. However, it does not capture localized or dynamic interventions. This limitation matters in Ethiopia in the study by Alemu et al \cite{Alemu2014}. There, malaria risk is clustered geographically. In such a setting, real control efforts would likely target high-risk villages rather than applied uniformly across the whole region.

\section{Modeling The Malaria Infectious System Around Disease-free Equilibrium}

\subsection{Linearization of The Malaria Infectious Model Around Disease-free Equilibrium}

From the preceding sections, we asserted that the population is conserved. This assumption allowed us to reach the non-dimensional system of PDEs that represents even mosquito control:
\begin{equation}
\[
\begin{cases}
\dfrac{\partial I_h}{\partial t}
=
D_h \dfrac{\partial^2 I_h}{\partial x^2}+\alpha I_m(1-I_h)-I_h, \\[6pt]
\dfrac{\partial I_m}{\partial t}
=
D_m\dfrac{\partial^2 I_m}{\partial x^2}
+
\dfrac{\delta}{\varepsilon}I_h(1-I_m)
-
\left(\dfrac{1}{\varepsilon}+c\right)I_m.
\end{cases}
\]

\end{equation}
To analyze the behavior of the malaria model near the disease free equilibrium, we assume that  near the disease-free equilibrium $(I_h, I_m) \approx (0,0)$. Under this assumption, we can apply Taylor linearization to argue that $I_mI_h \approx 0$, leading to the linearization of terms: 

\[
\alpha[I_m(1-I_h)]
=
\alpha[I_m-I_hI_m]
\approx
\alpha I_m,
\]
\[
\frac{\delta}{\varepsilon}[I_h(1-I_m)]
=
\frac{\delta}{\varepsilon}[I_h-I_hI_m]
\approx
\frac{\delta}{\varepsilon}I_h,
\]
\[
-\left(\frac{1}{\varepsilon}+c\right)I_m
\]
Applying this above approximations to the non-linear PDE system yields a linearized system:
\begin{equation}
\[
\begin{cases}
\dfrac{\partial I_h}{\partial \tau}
=
D_h\dfrac{\partial^2 I_h}{\partial x^2}
+
\alpha I_m
-
I_h, \\[6pt]
\dfrac{\partial I_m}{\partial \tau}
=
D_m\dfrac{\partial^2 I_m}{\partial x^2}
+
\dfrac{\delta}{\varepsilon}I_h
-
\left(\dfrac{1}{\varepsilon}+c\right)I_m.
\end{cases}
\]
\end{equation}
which can be expressed as a matrix equation as:
\begin{equation}
\label{eq:vector}
\frac{\partial \mathbf{U}}{\partial \tau}
= \mathbf{D}\,\mathbf{U}_{xx} + \mathbf{A}\,\mathbf{U}, \qquad
\mathbf{U} = \begin{pmatrix} I_h \\ I_m \end{pmatrix},
\end{equation}
with diffusion and reaction matrices
\begin{align}
\mathbf{D} = \begin{pmatrix} D_h & 0 \\ 0 & D_m \end{pmatrix},
\qquad
\mathbf{A} = \begin{pmatrix} -1 & \alpha \\ \delta/\varepsilon &
-(1/\varepsilon  + c)\end{pmatrix}.
\end{align}
It is important to note that the above linearization applies to the model for small perturbations in $I_h$ and $I_m$ around the disease-free equilibrium. Therefore, once infection levels become substantial, i.e $I_h, I_m \geq 0.2$, the original non-linear terms modify the linearized system, breaking every conclusion from the analysis and numerical computation around the linearized PDE system.

\subsection{Reduction of Linearized Malaria Model to the Malaria ODE system}

We achieve stability of the linearized malaria system when all eigenvalues of the system are negative, demonstrating the decay in infection over time from both the diffusion and reaction effects. To analyze this stability, we need to isolate the effects of diffusion from the reaction effects. \\  
Focusing on the diffusion effect, we solve a Sturm-Liouville problem generated by the diffusion, i.e $\Delta$ operator  under the imposed Neumann boundary conditions where the eigenvalues and eigenfunctions are expressed as:
\begin{equation}
\phi_n(x) = \cos\left(\frac{n\pi x}{L}\right),\quad \lambda_n = -\left(\frac{n\pi}{L}\right)^{2},\quad n = 0,1,2,\ldots
\end{equation}

Thus , we can rewrite $\mathbf{U}(x, \tau)$  as an infinite sum of the eigenfunctions:
\begin{equation}
\mathbf{U}(x,\tau) = \sum_{n=0}^\infty \mathbf{a}_n(\tau)\phi_n(x) \qquad
where \quad \mathbf{a}_n(\tau) = \begin{pmatrix}
a_{h,n}(\tau)\\ a_{m,n}(\tau)
\end{pmatrix}
\end{equation}

According to the Sturm-Liouville theory, the eigenfunctions of the Laplacian operator forms an orthogonal basis of the eigenfunctions' solution space\cite{Logan2015}.  Formally, the orthogonality argument can be represented as:
\begin{equation}
\int_0^L \phi_n(x)\,\phi_s(x)\,dx 
= \|\phi_s\|^2\, \delta_{n, m},
\end{equation} where         $\delta_{n, m}$ $=  \begin{cases}
    1 \quad n = m \\
    0  \quad elsewhere
\end{cases}

Substituting the eigenfunction expansion  $\mathbf{U}(x,\tau) = \sum_{n=0}^\infty 
\mathbf{a}_n(\tau)\phi_n(x)$ into matrix equation and using the property $\phi_n'' = \lambda_n 
\phi_n$, we observe that:
\begin{equation}
\sum_{n=0}^\infty \mathbf{a}_n'(\tau)\phi_n(x)
= \sum_{n=0}^\infty 
(\mathbf{A} + \lambda_n \mathbf{D})\,
\mathbf{a}_n(\tau)\phi_n(x).
\end{equation}
To isolate the mode m, we multiply both sides with $\phi_s$ and integrate over the region [0, L] which results into:
\begin{equation}
\sum_{n=0}^\infty \mathbf{a}_n'(\tau)
\int_0^L \phi_n(x)\phi_s(x)\,dx
=
\sum_{n=0}^\infty 
(\mathbf{A} + \lambda_n\mathbf{D})\,\mathbf{a}_n(\tau)
\int_0^L \phi_n(x)\phi_s(x)\,dx.
\end{equation}

Now applying the orthogonality argument, we observe that only the $n = m$ term survives on both sides, leading to a term-wise expression:
\begin{equation}
\mathbf{a}_s'(\tau)\|\phi_s\|^2
= (\mathbf{A} + \lambda_s\mathbf{D})\,
\mathbf{a}_m(\tau)\|\phi_s\|^2.
\end{equation}

Dividing both sides by $\|\phi_s\|^2 > 0$:
\begin{equation}
\mathbf{a}_n'(\tau) 
= \mathbf{M}_n\,\mathbf{a}_n(\tau), 
\qquad where \quad 
\mathbf{M}_n = \mathbf{A} + \lambda_n\mathbf{D}
= \begin{pmatrix}
-1 + \lambda_n D_h & \alpha \\
\frac{\delta}{\varepsilon }& 
-(\frac{1}{\varepsilon} +c) + {\lambda_n D_m}
\end{pmatrix}.
\end{equation}
Having reduced the linearized PDE system into an ordinary differential equation whose stability can be observed from its eigenvalues. Looking at the standard matrix, $M_m$, we observe that both its $tr(M_m) < 0 $ while $det(M_m) > 0 \quad  \forall \lambda_n < 0 \quad and \quad \delta, \epsilon, D_h, \ and\,D_m > 0,$ which are enforced by the biological constraints.

Therefore, we observe that the linearized model is reduced to the ODE model, which achieves stability in the localized region of infection levels around the disease-free equilibrium.

\subsection{Analysis of The Behaviour of The Malaria Infectious ODE Model }

By further analyzing the behavior of $M_s$ in the ODE model of the preceding section, we can discern the effects of diffusion and reaction in the spreading of the malaria in both human and mosquito population fraction, i.e $I_h, I_m$. 

Firstly, when $s = 0$, we observe that $\lambda_0 = 0, \  \phi_0 = 1.$ Thus, $M_0 = A$, which governs the reaction component of the ODE model. If the reaction matrix $A$ has negative eigenvalues, satisfying $\mathcal{R}_0 < 1$ where $\mathcal{R}_0 = \sqrt\frac{\alpha\delta}{1+c\epsilon}$ being the basic reproduction number, the system remains linearly stable, and there is no outbreak.  On the other hand, if $\mathcal{R}_0 > 1$, we observe an outbreak growth.

For $n \geq 1,$  $\lambda_s <0,$ we observe the stabilizing or damping effect of diffusion, represented by the spatial component of $M_s$ through the term $\lambda_n D. $ As $s$ increases, the dominant eigenvalue of $M_s$ becomes increasingly negative, resulting in high-frequency outbreak spikes to be damped faster across the spatial domain, despite the state of $\mathcal{R}_0$. Thus, this diffusing mechanism dampens the effect of the reaction mechanism, i.e A.

\section{Numerical Computation of Solutions and Stability Analysis}

\subsection{Finite Differences Scheme}
In real-world applications, solving the PDE system through infinite sums or similar methods
can be a very lengthy and computationally expensive approach to execute. Therefore, we
introduce the finite differences method, which leverages Taylor series and Euler’s method
to approximate the time derivative and spatial second derivatives, allowing the approxi-
mation of solutions across multiple temporal and spatial pairs for the complex non-linear
non-dimensional PDE system.

First, we discretize the spatial and temporal variables, i.e x and t, from their continuous
states.

Let \(x_j = jh\), with \(j = 0, 1, 2, ..., J\) and \(h = \frac{L}{J}\), representing one spatial step, for
\(x_j \in [0, L]\). Similarly, let \(t = nk\) with \(n = 0, 1, 2, . . . , N\) and \(k = \frac{T}{N}\), representing
one temporal step, for \(t_n \in [0, \tau]\).

By Taylor’s theorem, expanding \(u\) about \((x_j , t_n)\) gives
\[
u(x_j , t_{n+1}) = u^n_j + k u_\tau (x_j , t_n) + O(k^2),
\]
which rearranges to the forward-difference approximation
\[
\frac{\partial u}{\partial \tau}\bigg|^n_j
=
\frac{u^{n+1}_j - u^n_j}{k} + O(k).
\]

For the spatial second derivative, Taylor expansion gives
\[
u(x_{j+1}, t_n) = u^n_j + h u_x + \frac{1}{2} h^2u_{xx} + \frac{1}{6}h^3u_{xxx} + O(h^4),
\]
\[
u(x_{j-1}, t_n) = u^n_j - h u_x + \frac{1}{2} h^2u_{xx} - \frac{1}{6}h^3u_{xxx} + O(h^4),
\]
where derivatives on the right are evaluated at \((x_j , t_n)\). Adding the two expansions, the
odd-order terms cancel:
\[
u^n_{j+1} + u^n_{j-1} = 2u^n_j + h^2 u_{xx}\bigg|^n_j + O(h^4).
\]

Solving for the second derivative gives the centered-difference approximation
\[
\frac{\partial^2u}{\partial x^2}\bigg|^n_j
=
\frac{u^n_{j-1} - 2u^n_j + u^n_{j+1}}{h^2} + O(h^2).
\]

Substituting these approximations into the controlled reaction-diffusion system from Section 5 and
rearranging gives the explicit update scheme
\[
I^{n+1}_{h,j}
=
I^n_{h,j}
+
r_h
\left(
I^n_{h,j-1} - 2I^n_{h,j} + I^n_{h,j+1}
\right)
+
k\left[
\alpha I^n_{m,j}
\left(1-I^n_{h,j}\right)
-
I^n_{h,j}
\right],
\]
\[
I^{n+1}_{m,j}
=
I^n_{m,j}
+
r_m
\left(
I^n_{m,j-1} - 2I^n_{m,j} + I^n_{m,j+1}
\right)
+
k\left[
\frac{\delta}{\varepsilon} I^n_{h,j}
\left(1-I^n_{m,j}\right)
-
\left(\frac{1}{\varepsilon}+c\right)I^n_{m,j}
\right],
\]
where \(r_h = D_hk/h^2\) and \(r_m = D_mk/h^2\). The no-flux boundary conditions are imposed by
setting \(I^n_{\cdot,0} = I^n_{\cdot,1}\) and \(I^n_{\cdot,J} = I^n_{\cdot,J-1}\) at each time step, which enforces a zero-gradient condition
at the endpoints.

From the two equations expressed above, we observe that \(I(x_j , t_{n+1})\) is determined by 3
preceding values, \(I^n_j\), \(I^n_{j-1}\), and \(I^n_{j+1}\), which are all calculated at \(t = nk\). This allows us to
use recursion and make a scheme that draws a lattice of points that forms a grid, where
each value of \(I\) at the next time step, i.e., \(t = (n + 1)k\), is influenced by its preceding values,
\(t = nk\).

\subsection{Von Neumann Stability Analysis of the Finite-Difference Scheme}

We now analyze the stability of the diffusion part of the finite-difference scheme. Von Neumann stability analysis checks whether numerical errors in Fourier modes grow uncontrollably over time. For the scalar heat equation discretized by forward Euler in time and centered differences in space, we write a Fourier mode as
\[
I_j^n=G^n e^{i\beta jh}.
\]
Substituting this into the diffusion update gives
\begin{align}
G &= 1+r\left(e^{i\beta h}-2+e^{-i\beta h}\right)\\
&=1-4r\sin^2\left(\frac{\beta h}{2}\right).
\end{align}
For stability of the diffusion discretization, we need \(|G|\leq 1\) for all wave numbers \(\beta\). This gives the standard explicit heat-equation stability condition
\[
r\leq \frac{1}{2}.
\]
Applying this condition separately to the human and mosquito diffusion terms gives
\begin{equation}
\label{eq:stability}
r_h=\frac{D_hk}{h^2}\leq \frac{1}{2},
\qquad
r_m=\frac{D_mk}{h^2}\leq \frac{1}{2}.
\end{equation}

This von Neumann calculation controls the diffusion part of the explicit scheme. It does not prove that the full nonlinear reaction-diffusion update preserves the biological bounds
\[
0\leq I_h,I_m\leq 1.
\]
That would require a separate positivity-preservation analysis involving both the diffusion and reaction terms. Therefore, the stability condition above should not be interpreted as a proof that the numerical solution remains biologically meaningful for all parameter choices.

\section{Numerical Solutions of the Malaria Outbreak Model}

Using the method of finite differences we explored in the previous section, we model the behavior of the malaria spreading in the presence of low diffusion, as analyzed in the preceding sections. From the Von Neumann analysis above, we are confident our numerical approximation scheme is stable.

To achieve the method of finite difference computationally, we implement schemes in Python that operate vectorized slicing on the grid of discretized spatial and temporal positions for the full nonlinear system. In the numerical simulations, we use the parameters described below following the biological and physical constraints discussed in the preceding sections.
\begin{table}[H]
\centering
\caption{Scheme's Simulation Parameters.}
\begin{tabular}{|l|l|l|l|}
\hline
Parameter & Value & Parameter & Value \\
\hline
$L$ & 50 & $\alpha$ & 3.0 \\\hline
$J$ & 300 & $\delta$ & 2.67 \\\hline
$h = L/J$ & $0.1\overline{6}$ & $\varepsilon$ & 0.67 \\\hline
$k$ & $5\times10^{-4}$ & $x_0$ & 12.5 \\\hline
$D_h$ & 0.05 & $A_h = A_m$ & 0.5 \\\hline
$D_m$ & 0.15 & $\mathcal{R}_0 = \sqrt{\alpha\delta}$ (without c)& $ 2.83$\\
\hline
 c& 12& $\mathcal{R}_0 = \sqrt{\frac{\alpha\delta}{1+c\epsilon}}$ (with c)&0.941\\ \hline
\end{tabular}
\end{table}
In the low diffusion simulations, $D_m$ and $D_h$ are represented by values 0.05 and 0.15 as shown above. L and J is within the bounds that prevents infinite simulations for the non-linear system. We model the condition with the absence of mosquito control, i.e $c = 0$ where we have an outbreak modeled by $\mathcal{R}_0 > 1$ and in the presence of mosquito control where $c = 12$ where $\mathcal{R}_0 <1.$

\subsection{Outbreak Dynamics Without Control}
We first simulate the no-control case ($c = 0$), starting from a Gaussian initial profile centered at $x_0 = 12.5$. Since $\mathcal{R}_0 \approx 2.83 > 1$, Section~9.3 predicts that the disease-free equilibrium is unstable and the outbreak should grow toward an endemic equilibrium. 

\begin{figure}[H]
\includegraphics[width =0.5\textheight]{plot1.png}
\caption{Mosquito outbreak for $n=0$ mode, representing the absence of diffusive smoothing.}
\end{figure}

The graph above illustrates the behavior of the system when $n = 0$, i.e the absence of diffusion.  Both the human and mosquito infection population begin in a localized initial position $x_0 = \tfrac{L}{4}$. As $\tau$ increases, the $n=0$ mode where  the M = A dominates, with $\mathcal{R}_0 \approx 2.83,$ leading to malaria outbreak in both human and mosquitoes  as observed in the graph $\tau = 40$.

\begin{figure}[H]
\includegraphics[width=0.35\textheight]{plot2.png}
\caption{Effect of Mosquito Diffusion, $D_m$, on dampening the malaria outbreak, leading to overall spread}
\end{figure}
The graph above models the role of diffusion in creating spatial damping of the outbreak. We observe that for a low diffusion coefficient for mosquitoes, i.e., $D_m = 0.05$, the outbreak is concentrated largely near $x_0$. On the other hand, for a slightly higher mosquitoes diffusion coefficient, i.e., $D_m = 0.50$, we observe a more evenly diffusive damping with the infection broadening across the whole x domain as explained in the preceding section. These patterns align with real-world malaria data.\\
A study in northwest Ethiopia found that over 75\% \cite{Alemu2014} of cases were concentrated in high-altitude villages rather than spread uniformly across all villages, indicating limited effective diffusion and persistence of localized hotspots.

\begin{figure}[h]
    \centering
    \includegraphics[width=0.45\linewidth]{plot3.png}
    \caption{3D Numerical simulation of the malaria spreading model in humans, i.e $I_h$, for $\tau \in [0, 20]$ with $\mathcal{R}_0 \approx 2.83$ and $D_m = 0.15$}
\end{figure}

The graph above models the numerically approximated evolution of the infected human fraction $I_h(x, \tau)$ starting from an initial position, $x_0 = 12.5$ under the low  human and mosquito diffusion state. We observe vertical growth, which explains the reaction matrix between the interaction of human infection and mosquito infection, while the widening of the flow over time confirms our hypothesis of the spatial damping of malaria infection driven by the diffusion of the mosquitoes and humans. It is important to observe that the diffusive damping of the malaria system is achieved in low diffusion rates, confirming that all forms of diffusion cause the outbreak spreading across the spatial region rather than high concentrated spikes.

\subsection{Effect of Mosquito Control}

We also simulate the full nonlinear PDE with  control rate $c$ and compare the resulting space-time evolution. From the analytical threshold $\mathcal{R}_0^c = \sqrt{\frac{\alpha\delta} {1 + c\varepsilon}} < 1$, the critical control rate is
\[
c^\ast = \frac{\alpha\delta - 1}{\varepsilon} = \frac{8.01 - 1}{0.67} \approx 10.46
\]
Above $c^\ast$, the disease-free equilibrium is linearly stable and outbreaks should die out; below $c^\ast$, infection persists. We compare $c = 0$ (no control, $\mathcal{R}_0^c \approx 2.83$) with $c = 12$ (above threshold, $\mathcal{R}_0^c \approx 0.94$). The choice $c = 12$ is around 15\% above $c^\ast$, so that the suppression can be visible in $\tau \in [0, 20]$.

\begin{figure}[H]
\centering
\includegraphics[width=\linewidth]{plot4.png}
\caption{Spatial profiles of $I_h$ (solid blue) and $I_m$ (dashed red) at $\tau = 0, 10, 25, 40$ with $c = 12 > c^\ast \approx 10.5$. The initial Gaussian outbreak decays rapidly toward the disease-free equilibrium.}
\end{figure}
The figure above shows snapshots of $I_h$ and $I_m$ at $\tau = 0, 10, 25, 40$ for $c = 12$. The same Gaussian initial condition that produced an endemic outbreak in Figure~\ref{fig:plot1} now decays rapidly, by $\tau = 10$ the peak has dropped from 0.5 to roughly 0.04, by $\tau = 25$ infection levels are barely distinguishable from 0, and by $\tau = 40$ the system has returned to the disease-free state. The mosquito profile $I_m$ decays  faster than $I_h$ because the additional removal term $-c\,I_m$ acts directly on the mosquito population.
\begin{figure}[H]
\centering
\includegraphics[width=0.7\linewidth]{plot5.png}
\caption{Space-time evolution of $I_h(x, \tau)$ with  control ($c = 12$), $\tau \in [0, 20]$. The initial Gaussian (visible at $\tau = 0$, $x \approx 12.5$) decays as the additional mosquito mortality drives $\mathcal{R}_0^c$ below 1.}
\end{figure}

The 3D view in figure above shows the contrast with the no-control case (Figure~\ref{fig:plot3}) as the same initial condition that produced a wave invading the entire domain now produced a brief one and then vanished. The simulation's maximum value of $I_h$ at $\tau = 20$ is below $10^{-3}$, which fits the exponential decay at the rate predicted by the dominant eigenvalue of $\mathbf{M}_0 = \mathbf{A}$ when $c > c^\ast$. Thus, $c^\ast$ derived from linear stability analysis correctly predicts the qualitative behavior of the full nonlinear system. 

\section{Discussion and Conclusion}

This paper presented the Ross ODE model, and extended it into a reaction-diffusion PDE framework to capture the spatial nature of malaria transmission. The Ross model provides a clean foundation, identifying $\mathcal{R}_0 = \sqrt{\alpha\delta}$ as the threshold determining whether malaria dies out or persists. However, the assumption of well-mixed populations prevents the ODE model from representing where infection occurs or how it spreads across space.

Extending to a reaction-diffusion PDE framework allows infection to be tracked as a spatial fractions that evolves across both space and time, with movement entering through diffusion. Diffusion spreads infection outwards from concentrated areas over time, and the reaction terms determine whether total infection actually grows or decays at each location. 

We confirmed with our stability analysis that $\mathcal{R}_0$ remains the key threshold determining outbreak behavior through the $n = 0$ spatial mode, while higher modes are damped by diffusion regardless of $\mathcal{R}_0$. Numerical simulations with $\mathcal{R}_0$ $\approx$ 2.83 confirmed this, showing both human and mosquito infected fractions growing and spreading outward from the initial outbreak location, with $D_m$ controlling how broadly infection dispersed. When mosquito control was applied ($c = 12$), the infectious mosquito population dropped and transmission weakened ($\mathcal{R}_0$ $\approx$ 0.94). With a strong enough control parameter, the outbreak can be suppressed entirely.

However, several limitations remain. The model is one-dimensional, diffusion coefficients are assumed constant across space and time, and mosquito control is modeled as spatially uniform. In practice, mosquito behavior depends on factors such as local ecology and season, and intervention efforts are often localized and temporary. Future work could allow model parameters to vary across space, extend the model to two dimensions, and make control strategies more realistic, which would all contribute to making the model a more practical tool for public health decision-making.


\newpage
\begin{thebibliography}{9}
\bibitem{Qu2023}
Z. Qu, D. Patterson, L. M. Childs, C. J. Edholm, J. Ponce, O. Prosper, and L. Zhao, “Modeling immunity to malaria with an age-structured PDE framework,” SIAM Journal on Applied Mathematics, vol. 83, no. 3, pp. 1098–1125, 2023. 
\bibitem{Ross1911}
Ronald Ross, \textit{The Prevention of Malaria}, John Murray, London, 1911.
\bibitem{Alemu2014}
K.~Alemu, A.~Worku, Y.~Berhane, and A.~Kumie,
``Spatiotemporal clusters of malaria cases at village level, northwest Ethiopia,''
\bibitem{Sandstede2026}
B. Sandstede, Lecture notes for Applied Partial Differential Equations, Brown University, Providence, RI, Spring 2026.
\bibitem{Logan2015}
J.~D. Logan, \textit{Applied Partial Differential Equations}, 3rd~ed. Springer, New York, 2015.
\textit{Malaria Journal}, vol.~13, p.~223, 2014.
DOI:~\texttt{10.1186/1475-2875-13-223}
\end{thebibliography}

\end{document}

